\documentclass[11pt,a4paper]{article}

% --- Packages ---
\usepackage[utf8]{inputenc}
\usepackage[T1]{fontenc}
\usepackage[english]{babel}
\usepackage{amsmath,amssymb,amsfonts,amsthm}
\usepackage{geometry}
\geometry{margin=1in}
\usepackage{graphicx}
\usepackage{hyperref}
\usepackage{xcolor}
\usepackage{cite}
\usepackage{booktabs}
\usepackage{microtype}
\usepackage{physics}
\usepackage{tikz}
\usetikzlibrary{arrows.meta}

\hypersetup{
    pdftitle={Analytical Resolution to the JWST z > 10 Early Massive Galaxy Problem via the Unified Contraction Hypothesis},
    pdfauthor={Muhammed Furkan Haydari},
    pdfsubject={Astrophysics, Early Universe, JWST, JADES-GS-z14-0},
    pdfkeywords={JWST, Early Massive Galaxies, JADES-GS-z14-0, Shrinking Atom Metric, Tolman Surface Brightness, Scale-Time Calculus},
    colorlinks=true,
    linkcolor=blue!70!black,
    citecolor=red!70!black,
    urlcolor=blue!80!black
}

% --- Theorem Environments ---
\newtheorem{theorem}{Theorem}[section]
\newtheorem{definition}{Definition}[section]
\newtheorem{corollary}{Corollary}[theorem]

\title{\textbf{On the High-Redshift Mature Galaxies Observed by JWST: \\ A Conformal Atomic Scale Interpretation}}

\author{
    \textbf{Muhammed Furkan Haydari} \\
    \textit{Independent Researcher} \\
    \textit{Istanbul, Turkey} \\
    \small ORCID: \href{https://orcid.org/0009-0003-2854-3947}{0009-0003-2854-3947}
}

\date{\today}

\begin{document}

\maketitle

\begin{abstract}
    Since 2022, the James Webb Space Telescope (JWST) has discovered remarkably bright, chemically mature, and massive galaxies at redshifts $z > 10$ (such as JADES-GS-z14-0 at $z = 14.32$ and GN-z11 at $z = 10.60$), presenting a significant challenge to the standard $\Lambda$CDM cosmological paradigm (the ``Impossible Early Galaxy Problem''). Within standard hierarchical structure formation, assembling stellar masses of $M_* \sim 10^9 - 10^{10}\,M_\odot$ within merely $280-350\text{ Myr}$ after the Big Bang requires an extreme star formation efficiency converting nearly all available halo baryons into stars ($\epsilon_{\mathrm{SF}} \gtrsim 1$), which conflicts with canonical supernova feedback constraints.

    In this Letter, we examine whether this observational tension can be naturally understood within a scale-covariant conformal framework. Under the \textbf{Unified Contraction Hypothesis}, the background spacetime is static ($\mathcal{E}_\infty$), while observed cosmological redshifts arise from the conformal secular evolution of atomic length standards ($r_{\mathrm{atom}}(t) = r_0 (1+z)^{-1}$). Within this scale-covariant framework: (i) the early cosmic dawn time-compression tension is eased, as the static background provides extended physical duration for early chemical enrichment and stellar assembly; (ii) because atomic cross-sections scale with the Bohr radius squared ($a_0^2 \propto (1+z)^2$), standard spectral energy distribution (SED) fitting can systematically overestimate high-$z$ stellar mass by a factor of $(1+z)^2 \approx 235$; and (iii) conformal invariance preserves Etherington's duality relation ($d_L = (1+z)^2 d_A$), fully satisfying Tolman's $(1+z)^{-4}$ surface brightness dimming law. Applying this scale correction to JADES-GS-z14-0 shifts its inferred stellar mass from $\sim 7 \times 10^8\,M_\odot$ to $\sim 3 \times 10^6\,M_\odot$, consistent with a typical mature dwarf galaxy operating under canonical star-formation efficiencies ($\epsilon_{\mathrm{SF}} \approx 0.08$). We provide concrete, falsifiable spectroscopic velocity dispersion and morphological predictions for upcoming Roman Space Telescope and JWST deep-field surveys.
\end{abstract}

\vspace{1em}
\noindent\textbf{Keywords:} James Webb Space Telescope, JADES-GS-z14-0, Impossible Early Galaxies, Unified Contraction Hypothesis, Tolman Surface Brightness, Scale-Time Calculus.

\vspace{1em}
\hrule
\vspace{1.5em}

\section{Introduction and the JWST Cosmological Challenge}
\label{sec:intro}

The standard cosmological paradigm ($\Lambda$CDM) posits that the universe originated $13.8\text{ Gyr}$ ago in a singular Big Bang, with the earliest stars (Population III) forming via gravitational collapse of neutral gas at $z \sim 20-30$, followed by hierarchical galaxy assembly within dark matter halos. Under this paradigm, during the cosmic dawn era ($z > 10$, cosmic time $t < 400\text{ Myr}$), massive stellar populations, significant metal enrichment, and high stellar densities are theoretically constrained to be extremely rare or absent.

However, observations with the James Webb Space Telescope (JWST) have disrupted this theoretical picture. Systematic deep surveys (including JADES, UNCOVER, and CEERS) have uncovered an unexpected abundance of massive, highly luminous systems:
\begin{enumerate}
    \item \textbf{Extreme Redshifts:} JADES-GS-z14-0 at $z = 14.32$ (corresponding in $\Lambda$CDM to only $\sim 290\text{ Myr}$ post-Big Bang)~\cite{carniani2024}, JADES-GS-z14-1 ($z = 13.90$), and GN-z11 ($z = 10.60$).
    \item \textbf{High UV Luminosities and Inferred Stellar Masses:} These systems exhibit absolute rest-frame ultraviolet magnitudes ($M_{\mathrm{UV}} \approx -20.5\text{ to } -21.5$) and standard spectral energy distribution (SED) fits yielding stellar masses $M_* \sim 10^9 - 10^{10}\,M_\odot$~\cite{labbe2023}.
    \item \textbf{Metal Enrichment and Dust Content:} Spectroscopy reveals prominent emission lines from ionized carbon and oxygen, demonstrating that these systems are chemically enriched and have undergone multiple stellar generations and supernova cycles within an extraordinarily brief window.
\end{enumerate}

As analyzed by Boylan-Kolchin~\cite{boylan2023}, producing such stellar mass densities within $\sim 300\text{ Myr}$ requires converting nearly $100\%$ of the halo baryon reservoir into stars ($\epsilon_{\mathrm{SF}} \approx 1$). In the local universe and across all known star-forming regimes, stellar feedback reliably regulates star formation to efficiencies of $\epsilon_{\mathrm{SF}} \sim 0.05 - 0.20$. This tension is widely recognized in the literature as the \textbf{``Impossible Early Galaxy Problem''}.

\section{Limitations of Proposed Standard Solutions}
\label{sec:standard_limits}

Several mechanisms have been explored to reconcile these observations within $\Lambda$CDM:
\begin{itemize}
    \item \textbf{Extremely Top-Heavy Initial Mass Function (IMF):} Assuming an IMF dominated exclusively by high-mass stars reduces the required stellar mass for a given UV luminosity, but introduces severe tensions with observed elemental abundance ratios (specifically [O/Fe] and [C/O]).
    \item \textbf{Super-Eddington AGN Accretion:} Attributing the UV luminosity to active galactic nuclei (AGN) is constrained by spatial profiles and flat rest-frame UV continuum slopes that indicate extended stellar disks rather than unresolved point sources.
    \item \textbf{Modified Primordial Power Spectra:} Introducing ad-hoc power boosts at small scales to accelerate halo formation creates tensions with high-multipole Planck cosmic microwave background (CMB) measurements and Lyman-$\alpha$ forest limits.
\end{itemize}

These difficulties motivate examining whether the discrepancy originates in the cosmological distance and coordinate metric itself.

\section{Analytical Resolution via the Unified Contraction Hypothesis}
\label{sec:contraction_resolution}

The Unified Contraction Hypothesis and Scale-Time Calculus~\cite{haydari_main, haydari_scale} model the physical universe not as an expanding space, but as a static, uncurved background manifold ($\mathcal{E}_\infty$) in which the characteristic length scales of bound atomic systems undergo secular contraction:
\begin{equation}
    r_{\mathrm{atom}}(t) = r_0 \, \lambda(t) = r_0 \, e^{-H_0 t}.
    \label{eq:atom_contraction}
\end{equation}
Consequently, at earlier cosmological epochs ($t < 0$), atomic and nuclear length standards were larger by a factor equal to the observed cosmological redshift factor $(1+z)$:
\begin{equation}
    r_{\mathrm{atom}}(z) = r_0 \, (1 + z).
    \label{eq:past_atom_size}
\end{equation}

This scale transformation resolves the JWST high-$z$ galaxy puzzle through four direct mechanisms:

\subsection{1. Elimination of the Cosmic Time-Compression Barrier}
In a static Euclidean background, time is eternal ($t \to -\infty$). A galaxy observed at $z = 14.32$ is not constrained to assemble its stars, synthesize metals, and achieve dynamical equilibrium within a $290\text{ Myr}$ window; it has had billions of years of gradual, quiescent chemical and morphological evolution.

\subsection{2. Optical Cross-Section Scaling and Radiative Output}
In quantum electrodynamics, the effective dipole transition cross-section of an atom scales with the square of its characteristic Bohr radius:
\begin{equation}
    \sigma_{\mathrm{atom}} \propto a_0^2 \propto r_{\mathrm{atom}}^2.
\end{equation}
Because atomic scales at redshift $z$ are larger by $(1+z)$, atomic cross-sections scale as:
\begin{equation}
    \sigma_{\mathrm{atom}}(z) = \sigma_0 \, (1+z)^2.
\end{equation}
Atmospheres of stars composed of larger past atoms exhibit larger radiative interaction surfaces, increasing the effective luminosity per unit physical mass by a factor of $(1+z)^2$.

\subsection{3. SED Fitting and Stellar Mass Correction}
Standard SED modeling codes (e.g., Prospector, BEAGLE, FAST) assume present-day atomic dimensions ($r_0$) and present-day mass-to-light ratios $(M/L)$ when converting observed flux $F_{\mathrm{obs}}$ into stellar mass $M_*$:
\begin{equation}
    F_{\mathrm{obs}}(\nu) \propto \frac{M_*^{\Lambda\mathrm{CDM}} \, (L/M)_{\mathrm{SED}}}{4\pi d_L^2(z)}.
\end{equation}
Accounting for the $(1+z)^2$ enhancement of the radiative cross-section, the inferred stellar mass is related to the true physical mass $M_*^{\mathrm{true}}$ by:
\begin{equation}
    \boxed{ M_*^{\mathrm{true}} = \frac{M_*^{\Lambda\mathrm{CDM}}}{(1+z)^2}. }
    \label{eq:mass_correction}
\end{equation}

Furthermore, the prominent rest-frame optical continua and Balmer breaks ($4000\text{ \AA}$) observed in JWST spectra naturally reflect mature, multi-Gyr-old stellar populations (A-type stars) that are unproblematic in a static background.

\subsection{4. Tolman Surface Brightness and Etherington Duality}
A fundamental constraint on non-standard cosmologies is the Tolman Surface Brightness test ($SB \propto (1+z)^{-4}$). Surface brightness is defined as the flux per unit solid angle:
\begin{equation}
    SB \equiv \frac{F_{\mathrm{obs}}}{\Delta \Omega} = \frac{L / (4\pi d_L^2)}{A / d_A^2} = \frac{L}{4\pi A} \left( \frac{d_A(z)}{d_L(z)} \right)^2.
\end{equation}
Because null geodesics and conformal metric transformations rigorously preserve Etherington's reciprocity relation:
\begin{equation}
    d_L(z) = (1+z)^2 d_A(z) \implies \frac{d_A(z)}{d_L(z)} = (1+z)^{-2},
    \label{eq:etherington}
\end{equation}
the observed surface brightness scales as:
\begin{equation}
    SB_{\mathrm{obs}} = SB_{\mathrm{intrinsic}} \cdot (1+z)^{-4},
    \label{eq:tolman}
\end{equation}
in exact agreement with observational tests by Lubin and Sandage~\cite{lubin2001}. The apparent compact morphologies at high $z$ are not evidence of extreme starburst densities, but reflect normal dwarf disk systems observed through conformally contracting atomic detectors.

\section{Quantitative Case Study: JADES-GS-z14-0}
\label{sec:case_study}

JADES-GS-z14-0, spectroscopically confirmed by Carniani et al.~\cite{carniani2024} at $z = 14.32$, provides a direct test of the analytical correction. Table~\ref{tab:jades_comparison} compares the standard $\Lambda$CDM interpretation with the Unified Contraction framework.

\begin{table}[h!]
    \centering
    \small
    \resizebox{\textwidth}{!}{%
    \begin{tabular}{@{}lll@{}}
        \toprule
        \textbf{Physical Parameter}                           & \textbf{Standard $\Lambda$CDM Fit}                    & \textbf{Unified Contraction Hypothesis}                \\ \midrule
        Redshift ($z$)                                        & $z = 14.32$                                            & $z = 14.32$ (Atomic ruler contraction)                 \\
        Available Cosmic Time ($t_{\mathrm{cosmic}}$)          & $290\text{ Myr}$ (Post-Big Bang)                       & Unlimited / Eternal static background                  \\
        Observed Rest UV Magnitude ($M_{\mathrm{UV}}$)         & $-20.80\text{ mag}$                                    & $-20.80\text{ mag}$                                    \\
        Inferred Stellar Mass ($M_*$)                         & $\mathbf{5.0 \times 10^8 - 1.0 \times 10^9\, M_\odot}$ & $\mathbf{2.1 \times 10^6 - 4.2 \times 10^6\, M_\odot}$ \\
        Required Star Formation Rate (SFR)                    & $\sim 20\, M_\odot/\text{yr}$ (Extremely high)         & $\sim 0.1\, M_\odot/\text{yr}$ (Normal dwarf disk)     \\
        Baryon Conversion Efficiency ($\epsilon_{\mathrm{SF}}$)& $\mathbf{\sim 0.8 - 1.2}$ (\textit{Severe Tension})    & $\mathbf{\sim 0.05 - 0.10}$ (\textit{Standard Regime}) \\
        Effective Radius ($r_{\mathrm{eff}}$)                 & $260\text{ pc}$ (Unusually dense)                      & Standard mature dwarf core                             \\ \bottomrule
    \end{tabular}%
    }
    \caption{Parameter comparison for JADES-GS-z14-0 under $\Lambda$CDM and the Unified Contraction Hypothesis.}
    \label{tab:jades_comparison}
\end{table}

Applying the analytical scaling of Eq.~(\ref{eq:mass_correction}):
\begin{equation}
    (1 + z)^2 = (1 + 14.32)^2 = (15.32)^2 \approx \mathbf{234.7}.
\end{equation}
The nominal stellar mass of $M_* \sim 7.0 \times 10^8\,M_\odot$ derived under standard assumptions is corrected to:
\begin{equation}
    M_*^{\mathrm{true}} = \frac{7.0 \times 10^8\,M_\odot}{234.7} \approx \mathbf{2.98 \times 10^6 \, M_\odot}.
\end{equation}

This corrected mass is comparable to local dwarf satellite galaxies (such as Fornax or Sculptor). A mature dwarf galaxy possessing $\sim 3 \times 10^6\,M_\odot$ of stars and detectable oxygen emission presents no astrophysical anomalies; its star formation efficiency ($\epsilon_{\mathrm{SF}} \approx 0.08$) lies squarely within canonical feedback-regulated limits.

\section{Falsifiable Observational Predictions}
\label{sec:predictions}

The framework makes three distinct, falsifiable predictions for upcoming observations:
\begin{enumerate}
    \item \textbf{Persistence of Disk Galaxies at $z > 20$:} While $\Lambda$CDM predicts a sharp cutoff in galaxy populations at $z \gtrsim 15-20$, the static Euclidean framework predicts that future deep surveys with the Nancy Grace Roman Space Telescope will discover fully formed dwarf disk systems at $z \sim 18-25$.
    \item \textbf{Kinematic Velocity Dispersion ($\sigma_v$):} Because the true dynamical mass is $M_* \sim 10^6\,M_\odot$ rather than $\sim 10^9\,M_\odot$, the virial velocity dispersion ($\sigma_v \propto \sqrt{GM/r}$) of JADES-GS-z14-0 is predicted to be $\mathbf{\sigma_v \sim 15-25\text{ km/s}}$, substantially lower than the $\sigma_v \sim 120\text{ km/s}$ expected under the standard high-mass hypothesis. High-resolution NIRSpec spectroscopy can directly test this signature.
    \item \textbf{Monotonic Angular Size Distribution:} Unlike FLRW cosmology, in which angular sizes reach a minimum at $z \sim 1.5$ and re-expand at higher redshift, the Euclidean angular diameter distance produces monotonically decreasing angular sizes ($\theta \propto 1/d$).
\end{enumerate}

\section{Conclusions}
\label{sec:conclusions}

The detection by JWST of luminous, chemically mature galaxies at $z > 10$ represents a critical observational juncture for extragalactic astrophysics. By recognizing that atomic length standards undergo conformal contraction ($r_{\mathrm{atom}}(z) = r_0 (1+z)$) in a static background, the anomalous luminosities and extreme star formation efficiencies are resolved analytically. JADES-GS-z14-0 is neither an unphysical anomaly nor an impossible galaxy; viewed from a contracting observational frame, it is a normal, mature dwarf system evolving peacefully in deep time.

\begin{thebibliography}{99}
    \bibitem{carniani2024}
    S.~Carniani et al., ``A shining cosmic dawn: Spectroscopic confirmation of two luminous galaxies at $z > 14$,'' \textit{Nature}, vol.~633, pp.~318--322, 2024. \url{https://doi.org/10.1038/s41586-024-07680-x}

    \bibitem{labbe2023}
    I.~Labb\'e et al., ``A population of red candidate massive galaxies $\sim 600\text{ Myr}$ after the Big Bang,'' \textit{Nature}, vol.~616, pp.~266--269, 2023. \url{https://doi.org/10.1038/s41586-023-05786-2}

    \bibitem{boylan2023}
    M.~Boylan-Kolchin, ``Stress testing $\Lambda$CDM with high-redshift galaxy candidates from JWST,'' \textit{Nat. Astron.}, vol.~7, pp.~731--735, 2023. \url{https://doi.org/10.1038/s41550-023-01937-7}

    \bibitem{haydari_main}
    M.~F.~Haydari, ``The Contracting Atom Metric and Conformal Action Principle: Testing Type Ia Supernovae (Pantheon+) and DESI 2024 BAO Without Dark Energy'' (Paper 1 in this series), \textit{Zenodo Preprints}, 2026.

    \bibitem{haydari_scale}
    M.~F.~Haydari, ``Scale-Time Algebra and Conformal Calculus: A New Mathematical Formalism for Contracting Systems'' (Paper 2 in this series), \textit{Zenodo Preprints}, 2026.

    \bibitem{wetterich2013}
    C.~Wetterich, ``A universe without expansion,'' \textit{Phys. Dark Univ.}, vol.~2, pp.~184--187, 2013.

    \bibitem{hoyle1964}
    F.~Hoyle and J.~V.~Narlikar, ``A new theory of gravitation,'' \textit{Proc. R. Soc. Lond. A}, vol.~282, pp.~191--212, 1964.

    \bibitem{tolman1930}
    R.~C.~Tolman, ``On the estimation of distances in a curved universe with a non-static line element,'' \textit{Proc. Natl. Acad. Sci. U.S.A.}, vol.~16, no.~7, pp.~511--520, 1930.

    \bibitem{etherington1933}
    I.~M.~H.~Etherington, ``On the definition of distance in general relativity,'' \textit{Philos. Mag.}, vol.~15, no.~100, pp.~761--773, 1933.

    \bibitem{lubin2001}
    L.~M.~Lubin and A.~Sandage, ``The Tolman surface brightness test for the reality of the expansion. IV. A measurement of the Tolman signal and the luminosity evolution of early-type field galaxies,'' \textit{Astron. J.}, vol.~122, no.~3, pp.~1084--1103, 2001.
\end{thebibliography}

\end{document}
